\section{Introduction}\label{sec:introduction}

A random curve whose dimension lies strictly between one and two
requires a finer notion of length than arc length. A dimension theorem
specifies the critical power of a small covering diameter, but that
power alone need not give a nonzero measure. For ordinary chordal
$\mathrm{SLE}_\kappa$, $0<\kappa<8$, Beffara identified the dimension
$d=1+\kappa/8$ \cite[main Theorem]{beffara2008}; Rezaei proved that its
$d$-dimensional Hausdorff measure is zero
\cite[Theorem~1.1]{rezaei2018}. The remaining question concerns the
correction to the critical power, not a further change of dimension.

A \emph{gauge} is a continuous nondecreasing function
$h:[0,\infty)\to[0,\infty)$ with $h(0)=0$ and $h(r)>0$ for $r>0$.
Its Hausdorff measure,
in the diameter convention used here, is
\begin{equation}\label{eq:hausdorff-definition}
 \mathcal H^h(E)=\lim_{\delta\downarrow0}\inf\left\{
 \sum_j h(\operatorname{diam}U_j):
 E\subset\bigcup_jU_j,\quad \operatorname{diam}U_j\le\delta\right\},
\end{equation}
where covers are countable. We write $\mathcal H^h_\delta$ for the
infimum at a fixed cutoff. We seek a deterministic gauge that gives
positive finite measure to each nontrivial compact trace segment.
Hausdorff measure optimizes over spatial covers; ordered partitions
of a parametrized curve lead to a different notion of length.

Schramm asked for a sigma-finite Hausdorff measure of the SLE trace
and suggested the following candidate
\cite[Problem~7.1 and the following discussion]{schramm2007}:
\[
 r^d\log\log(1/r).
\]
The article \emph{An exact Hausdorff gauge for SLE: A moment-integral
and finite-batch construction} already answered the existence question
for each fixed $0<\kappa<8$. Its deterministic moment integrals specify
coefficients, and finite irregular batches of radii specify a gauge
through a full infimum \cite[equations~(5.8) and~(6.4)]{openai2026exact}.
That gauge gives positive finite measure simultaneously on positive-time
compact segments and finite expected measure of the entire trace in
bounded regions \cite[Theorem~1.1 and Proposition~8.2]{openai2026exact}.

The present paper identifies a regular closed-form gauge and proves
its two Hausdorff inequalities from quantitative estimates at finite
scales. The stopped-law, mass, and covering arguments below do not
import a proof result from the companion. The relation between the
two exact gauges remains open.
The regular formula also settles the proposed exponent-one correction
negatively: its measure is not sigma-finite on any positive-time
compact segment.

\subsection{The result}

Fix $0<\kappa<8$ and set
\begin{equation}\label{eq:parameters}
 d=1+\frac\kappa8,\qquad p=2-d=1-\frac\kappa8.
\end{equation}
Thus $1<d<2$ and $0<p<1$. Define, at all sufficiently small radii,
\begin{equation}\label{eq:gauge}
 h(r)=r^d\bigl(\log\log(1/r)\bigr)^{p/2}.
\end{equation}
Set $h(0)=0$ and complete it continuously and nondecreasingly at
larger radii. The expression is increasing near zero, and its
completion does not affect \eqref{eq:hausdorff-definition}.

\begin{theorem}\label{thm:main}
Let $\gamma$ be ordinary chordal $\mathrm{SLE}_\kappa$ from $0$ to
$\infty$ in $\mathbb H$, parametrized by half-plane capacity $2t$,
and set $\Gamma=\gamma([0,\infty))$. For each fixed $0<\kappa<8$,
on one event of probability one,
\[
 0<\mathcal H^h(\gamma([s,t]))<\infty
 \qquad\text{for every real }0<s<t<\infty.
\]
For every deterministic $0<R<\infty$,
\[
 \mathbb E\mathcal H^h(\Gamma\cap\overline B(0,R))
       \le C_{\kappa,R}<\infty.
\]
In particular, on the same probability-one event,
\[
 \mathcal H^h\bigl(\Gamma\cap([-m,m]+i[0,m])\bigr)<\infty
 \qquad(m=1,2,\ldots).
\]
\end{theorem}

The spatial expectation concerns the whole trace at once, including
its real-axis points. It is not obtained by adding bounds over
successive finite time intervals. The event is for the fixed value
of $\kappa$.

\begin{corollary}\label{cor:larger-gauge}
Let $g(r)=r^d\log\log(1/r)$ at sufficiently small radii, with any
continuous nondecreasing completion. On the event in
Theorem~\ref{thm:main}, the restriction of $\mathcal H^g$ to every
$\gamma([s,t])$, $0<s<t<\infty$, is not sigma-finite; in particular
its total measure is infinite.
\end{corollary}
\begin{proof}
The ratio $h(r)/g(r)$ tends to zero. For each $\varepsilon>0$,
all sufficiently fine covers therefore give
$\mathcal H^h(A)\le\varepsilon\mathcal H^g(A)$.
Every set of finite $\mathcal H^g$-measure has zero
$\mathcal H^h$-measure. A countable cover of the segment by such sets
would contradict its positive $\mathcal H^h$-measure.
\end{proof}

Natural parametrization provides another notion of SLE length.
Lawler--Sheffield constructed it for
$0<\kappa<4(7-\sqrt{33})$ \cite[Theorem~3.1]{lawlersheffield2011};
Lawler--Zhou extended the construction to $0<\kappa<8$
\cite[Theorem~1]{lawlerzhou2013}; and Lawler--Rezaei identified it,
up to normalization, with $d$-dimensional Minkowski content
\cite[Theorem~1.1 and p.~5]{lawlerrezaei2015}.
For chordal SLE in the half-plane, Holden--Yuan identify positive
fine-mesh variation, up to a deterministic factor, with natural
parametrization, including segments that meet the real boundary
\cite[Theorems~1.7--1.8]{holdenyuan2026}.
Their gauge $r^d(\log\log(1/r))^{-(d-1)}$ is applied to ordered
partitions. The different logarithmic correction here belongs to
arbitrary spatial covers. Our argument requires no identification
of $\mathcal H^h$ with either natural parametrization or Minkowski
content.

\subsection{The two Hausdorff bounds}

For the lower Hausdorff bound we need a positive measure on a piece
of the trace whose mass in every small ball is at most a constant
times $h(r)$. For the upper bound we need spatial covers whose total
$h$-cost has a bounded expectation. Both tasks use finite-mesh
quantities, but they use them in different ways.

First consider a disk $A$ with a fixed free collar in the surviving
domain after a finite stopped past. This domain is the component
toward infinity after filling the components cut off by the trace.
At a live interior point $z$, the Green density is
$R^{-p}\sin^{8/\kappa-1}\theta$, where $R$ is its conformal radius
and $\theta$ is its argument in tip--target half-plane coordinates.
Stop this density at the point's own first live approach within
distance $e$, and assign zero if no such approach occurs. Denote
the resulting stopped weight by $W_e(z)$. The quantitative estimate is
\[
 \mathbb P\left\{\int_AW_e(z)\,dA(z)>u\right\}
        \le C\exp(-cu^{2/p})
\]
in unit-scale coordinates, for large $u$, uniformly in the permitted
past and final mesh. This controls a fixed spatial scale while the
approximation mesh decreases to zero. Rezaei and Zhan proved
fixed-$n$ multipoint approach bounds and finiteness of moments of natural
length \cite[Theorems~1.1--1.2]{rezaeizhan2017}. Their estimates already
allow arbitrary approach radii. The estimate here gives the quantitative
$e^{-cu^{2/p}}$ tail for the integrated stopped density, with constants
uniform over both the permitted pasts and the final meshes.

The exponent comes from testing a large mass first at an intermediate
distance $\delta$. There the deterministic density bound gives mass
of order $\delta^{-p}$. For each prescribed finer final mesh, the
later increments raise this bound by only a fixed factor outside an
event of probability at most $Ce^{-c\delta^{-2}}$, with constants
independent of that final mesh. Taking $\delta$ of order $u^{-1/p}$,
and using the pointwise cap for coarser meshes, produces the displayed
tail. Section~\ref{mesh:section} controls those increments
in two ways. For $\kappa\le4$, Green-kernel decrease bounds the bracket
of density increments that start at different times. For $\kappa>4$,
spatially separated windows and a swallowing estimate control the
surviving contributions.

For the lower bound, Section~\ref{retain:section} deletes dyadic
squares whose mass is too large. At scale $r=2^{-k}$ the threshold is
$Dr^d(\log k)^{p/2}$ for a sufficiently large fixed constant $D$.
Reaching the square costs $O(r^p)$, whereas
its mass scale is $r^d$ and there are $O(r^{-2})$ squares.
The powers cancel because $p+d=2$. The tail makes the remaining
loss summable. Positive first moments and bounded second moments
leave a positive retained mass with positive probability. A supported
weak subsequential limit satisfies the required bound on every ball;
restarting on deterministic capacity slots gives positivity on every
positive-time interval.

The upper bound asks for the opposite kind of local event: a disk
must occasionally contain \emph{more} sausage area than its usual
scale. For a curve segment $\eta$, its normalized sausage area at
mesh $v$ is
\[
 v^{-p}\operatorname{Area}\{z:\operatorname{dist}(z,\eta)<v\}.
\]
Rezaei's covering proof already combines dense occupation, conditional
estimates after stopped histories, compactness, and prescribed spiral
passages \cite[Section~3, Lemmas~3.1--3.3]{rezaei2018}. Using the same
geometric idea, a serpentine ribbon with order $m$ strands provides order $m^2$
disjoint local tests at spatial scale $m^{-1}$. Successful tests
contribute total normalized area of order $m^{2-d}=m^p$. If each
step succeeds with a conditional probability bounded below uniformly
over every admissible rough past, the whole visit costs at most an
exponential in $m^2$: its probability is at least $ce^{-Cm^2}$.
Pointwise support for a fixed smooth path, as in Tran--Yuan
\cite[Theorem~1.1, Corollary~1.2, and Proposition~1.4]{tranyuan2020},
does not supply this common conditional probability over rough pasts.

Section~\ref{retain:section} first proves an ordinary-law lower-tail
statement for sausages on deterministic time intervals. In
Section~\ref{rc:section}, compactness of conformal maps along an interior
guide and a uniform localization estimate turn pointwise support into
one ribbon-step bound. The passage is described with fixed margins that
persist under convergence of the maps. A finite selection of rational
time intervals then transfers the sausage estimate, leaving a common
positive probability, mass threshold, and deterministic mesh cutoff.

For the final cover take $m$ proportional to $\sqrt{\log n}$ and
perform order $n$ tests at geometrically spaced radii $r_j$, with $j$ between
$n/3$ and $2n/3$. Conditional success bounds, not independence, make
total failure small. Since $\log n$ is comparable to
$\log\log(1/r_j)$, the generated mass matches the correction in
$h$. The deterministic mesh cutoff in the ribbon estimate lets
every test use one physical mesh. Disjoint selected disks can then
all be charged to the same whole-trace sausage measure.

Section~\ref{sec:dense-cover} carries out this cover. Its dense-occupation
and multiscale strategy follows Rezaei \cite[Section~3, especially
equations~(9)--(11)]{rezaei2018}, who selects maximal high-occupation
circles and charges them using bounded overlap. Here a greedy family
of disjoint disks is charged to one finite-mesh sausage measure, and
the quantitative ribbon probability specifies the logarithmic gauge.
To bound the cost of visited cells that the selected disks miss, the
proof uses test failures visible at finite stops, where the finite Green
changes of law apply. The disks themselves are selected afterward from
the whole-trace sausage. The final charge has a bounded expectation in
one fixed spatial disk, independently of time and height cutoffs.
Increasing both cutoffs then yields the all-time assertion.

The technical foundations appear in Sections~\ref{geom:section}
and~\ref{ll:section}. Starting from Rohde--Schramm's continuous trace
theorem \cite[Theorem~5.1]{rohdeschramm2005} and classical conformal
boundary theory, the first establishes access and slit support after a
rough past. The second develops the finite stopped Green laws from
the methods of Lawler--Zhou \cite[Section~2]{lawlerzhou2013} and
Lawler--Werness \cite[Sections~2.3--2.4]{lawlerwerness2013},
and proves the physical localization used by the ribbon argument.
The $\kappa\le4$ mesh proof uses the Green-decrease identity appearing in
Schramm--Sheffield's SLE$_4$/Gaussian-free-field covariance calculation
\cite[Lemmas~2.5--2.6]{schrammsheffield2013}, with the decomposition for
different starting times supplied here.
